\documentclass[conference]{IEEEtran}
\usepackage{cite}
\usepackage{amsmath,amssymb,amsfonts}
\usepackage{algorithmic}
\usepackage{graphicx}
\usepackage{textcomp}
\usepackage{xcolor}

\begin{document}

\title{Performance Analysis of a 3D Lennard-Jones Molecular Dynamics Simulator
       using OpenMP and MPI Parallelisation}

\author{\IEEEauthorblockN{HPC Student}
\IEEEauthorblockA{\textit{Department of Computer Science} \\
\textit{IIT Mandi}\\
Mandi, India \\
user@iitmandi.ac.in}}

\maketitle

\begin{abstract}
This paper presents the design, implementation, and performance evaluation of a 3D
Lennard-Jones Molecular Dynamics (MD) simulator in C++.
Starting from a verified serial baseline, we develop shared-memory (OpenMP) and
distributed-memory (MPI) parallel implementations, followed by a hybrid MPI+OpenMP
configuration.
Verification confirms NVE energy conservation ($<$0.61\% drift, $N=256$, $\Delta t=0.005$),
correct thermostat-driven temperature equilibration to $T_\mathrm{target}=2.0$ (LJ units),
and accurate radial distribution functions.
Timer-based profiling reveals that \texttt{compute\_forces} consumes $>$99.3\% of
wall-clock time for all $N$, motivating parallelisation.
Serial complexity scaling confirms $T \propto N^{1.70}$, approaching $O(N^2)$.
The OpenMP implementation achieves $5.90\times$ speedup at 8 threads (73.7\% efficiency)
relative to the serial baseline, with a fitted Amdahl serial fraction of
$f_\mathrm{serial}\approx5.2\%$, implying $S_\mathrm{max}\approx19\times$.
The MPI domain-decomposition implementation achieves $6.94\times$ speedup at $P=8$ ranks
(86.8\% efficiency) with $>$81\% weak-scaling efficiency at $P=8$.
The hybrid $4\text{r}\times4\text{t}$ configuration yields the best throughput
($13.2\times$ over serial).
\end{abstract}

\section{Introduction}
Molecular Dynamics (MD) is a cornerstone of computational physics and materials
science~\cite{b1,b2}, enabling atomic-level simulation of phase transitions, protein
folding, and materials properties.
The quadratic scaling of all-pair force evaluation renders MD a canonical HPC
target: for $N$ particles, every time step requires evaluating $N(N-1)/2$
pair interactions under a cutoff radius $r_c$.

This paper documents a complete MD pipeline beginning with a modular, cache-friendly
serial implementation in C++ and extending to OpenMP~\cite{b3}, MPI~\cite{b4},
and hybrid parallel strategies.
Each stage includes verification tests and quantitative scaling studies.
All timing uses a high-resolution \texttt{std::chrono} harness for consistency.

\section{Physical Model}
The system simulates $N$ particles in a cubic periodic box of side $L$,
interacting via the Lennard-Jones (LJ) potential:
\begin{equation}
  V(r) = 4\epsilon \!\left[ \!\left(\frac{\sigma}{r}\right)^{12}
                           \!-\! \left(\frac{\sigma}{r}\right)^{6} \right],
  \label{eq:lj}
\end{equation}
with cutoff $r_c = 2.5\sigma$ and minimum-image PBC.
Reduced LJ units ($\epsilon=\sigma=m=1$) are used throughout.
Pressure is computed via the virial theorem:
$P = \rho k_BT + \frac{1}{3V}\sum_{i<j}\mathbf{r}_{ij}\cdot\mathbf{F}_{ij}$.
At $\rho=0.8$, $T=1.0$, the simulation yields $P\approx2.27$, within 3.2\% of the
van~der~Waals equation-of-state prediction $P\approx2.20$, confirming thermodynamic
correctness.

\section{Numerical Algorithm}
Time integration uses the velocity Verlet algorithm~\cite{b7}:
\begin{align}
  \mathbf{v}(t + \tfrac{\Delta t}{2}) &= \mathbf{v}(t) + \tfrac{\mathbf{F}(t)}{m}\tfrac{\Delta t}{2}, \\
  \mathbf{r}(t + \Delta t)            &= \mathbf{r}(t) + \mathbf{v}(t + \tfrac{\Delta t}{2})\Delta t, \\
  \mathbf{v}(t + \Delta t)            &= \mathbf{v}(t+\tfrac{\Delta t}{2}) +
                                          \tfrac{\mathbf{F}(t+\Delta t)}{m}\tfrac{\Delta t}{2}.
\end{align}
This symplectic integrator conserves phase-space volume exactly,
giving $O(\Delta t^2)$ global truncation error and excellent long-time energy conservation.

\section{Serial Implementation}
Particle coordinates and forces use a Structure-of-Arrays (SoA) layout,
enabling SIMD auto-vectorisation.
Fig.~\ref{fig:compiler} confirms \texttt{-O3} provides a $4.2\times$ speedup over
\texttt{-O0} through loop unrolling and vectorisation of the LJ kernel.

\section{Verification and Testing}

\subsection{Energy Conservation (NVE)}
\begin{figure}[htbp]
\centerline{\includegraphics[width=\columnwidth]{plots/energy_conservation.pdf}}
\caption{NVE energy conservation: KE/N (orange, left axis) and PE/N, $E_\mathrm{tot}$/N
(blue/green, right axis). Inset: total energy drift as a percentage of the mean.
$(N=256,\;\Delta t=0.005,\;T_\mathrm{target}=2.0)$.}
\label{fig:nve}
\end{figure}
Fig.~\ref{fig:nve} presents per-particle kinetic, potential, and total energy
over 50 LJ time units on a dual y-axis, making both KE thermal equilibration and
PE relaxation simultaneously visible.
Total energy drift is 0.61\%, well within acceptable limits for $\Delta t=0.005$.

\subsection{Temperature Equilibration}
\begin{figure}[htbp]
\centerline{\includegraphics[width=\columnwidth]{plots/temperature_equilibration.pdf}}
\caption{Temperature equilibration: velocity-rescaling thermostat (blue, $t\le10$)
drives the system to $T_\mathrm{target}=2.0$; NVE production (orange, $t>10$) maintains
$\langle T\rangle\approx1.12$. Dashed: $T_\mathrm{target}=2.0$.}
\label{fig:equil}
\end{figure}
Fig.~\ref{fig:equil} shows the two-phase protocol: a velocity-rescaling thermostat is
applied for $t\le10$ (2000 steps), then removed for NVE production.
The production temperature $\langle T\rangle\approx1.12$ is lower than $T_\mathrm{target}=2.0$
because the velocity-rescaling thermostat sets the kinetic temperature, but upon removal,
potential energy stored in the compressed FCC lattice is released, converting kinetic
to potential energy and reducing the instantaneous temperature.
This is physically correct NVE behaviour: total energy is conserved while the
kinetic/potential partition adjusts to the true equilibrium.

\subsection{Radial Distribution Function and Convergence}
\begin{figure}[htbp]
\centerline{\includegraphics[width=\columnwidth]{plots/rdf.pdf}}
\caption{Radial Distribution Function $g(r)$ at $\rho=0.8$, $T=1.0$,
showing the expected first peak at $r\approx1.12\sigma$ and liquid-state oscillations.}
\label{fig:rdf}
\end{figure}
\begin{figure}[htbp]
\centerline{\includegraphics[width=\columnwidth]{plots/timestep_convergence.pdf}}
\caption{Velocity Verlet convergence: $|x-x_\mathrm{exact}|$ vs $\Delta t$ for a
harmonic oscillator confirms $O(\Delta t^2)$ convergence.}
\label{fig:verlet}
\end{figure}
Fig.~\ref{fig:rdf} confirms liquid structure at $\rho=0.8$, $T=1.0$.
Fig.~\ref{fig:verlet} confirms second-order convergence against the analytic solution.

\section{Code Profiling and Complexity}

\begin{table}[htbp]
\caption{Serial Profiling Results: Wall Time (s) by Function\newline
{\footnotesize 500 steps per run, $\rho=0.8$, $r_c=2.5$}}
\label{tab:profiling}
\begin{center}
\begin{tabular}{|l|c|c|c|c|c|}
\hline
\textbf{Function} & \textbf{N=100} & \textbf{N=200} & \textbf{N=500} & \textbf{N=1000} & \textbf{N=2000} \\
\hline
\texttt{compute\_forces}  & 0.0313 & 0.0972 & 0.4505 & 1.4296 & 5.4207 \\
\texttt{velocity\_verlet} & 0.0004 & 0.0007 & 0.0016 & 0.0032 & 0.0065 \\
\texttt{apply\_pbc}       & 0.0001 & 0.0003 & 0.0006 & 0.0011 & 0.0026 \\
I/O \& Thermo             & 0.0001 & 0.0001 & 0.0002 & 0.0003 & 0.0002 \\
\hline
TOTAL                     & 0.0319 & 0.0982 & 0.4529 & 1.4342 & 5.4300 \\
\hline
\end{tabular}
\end{center}
\vspace{2pt}
{\footnotesize All timings collected with \texttt{std::chrono} high-resolution clock on a
single core. \texttt{compute\_forces} consumes $>$99.3\% of total time at $N=1000$.}
\end{table}

Table~\ref{tab:profiling} quantifies the bottleneck:
\texttt{compute\_forces} dominates at $>$99.3\% of wall time for $N=1000$,
with the fraction growing with $N$ as the $O(N^2)$ pair loop dominates.
The $O(N)$ \texttt{velocity\_verlet} and \texttt{apply\_pbc} remain $<$1\% at all sizes.

\begin{figure}[htbp]
\centerline{\includegraphics[width=\columnwidth]{plots/md/serial_n2_scaling.pdf}}
\caption{Serial wall time vs $N$ on log-log axes (100 steps each).
Fitted $T\propto N^{1.70}$ ($R^2=0.9997$). The exponent approaches 2.0 at large $N$;
at moderate $N$, the cutoff sphere reduces the effective interaction count below $N^2/2$.}
\label{fig:n2scale}
\end{figure}

Fig.~\ref{fig:n2scale} confirms super-linear (approaching $O(N^2)$) growth.
The fitted exponent $\alpha=1.70$ is physically correct: at moderate $N$,
a significant fraction of pairs lies within $r_c$, reducing the effective
exponent; as $N\to\infty$, $\alpha\to2.0$.

\begin{figure}[htbp]
\centerline{\includegraphics[width=\columnwidth]{plots/compiler_optimisation.pdf}}
\caption{Wall-time speedup vs GCC optimisation level relative to \texttt{-O0}.
\texttt{-O3} achieves $4.2\times$ via SIMD vectorisation and loop unrolling.}
\label{fig:compiler}
\end{figure}

\section{OpenMP Parallelisation}

The outer loop of \texttt{compute\_forces} iterates over particle pairs $(i,j)$
with $i<j$.
Newton's third law allows accumulating $\mathbf{F}_{ji}=-\mathbf{F}_{ij}$ within
the same iteration, but naïve parallelisation of the $i$-loop creates a data race on
the force arrays.
We resolve this with \emph{thread-private} force arrays (one copy of size $N$ per thread)
accumulated into the global array inside a \texttt{\#pragma omp critical} section.
This avoids atomic operations in the hot inner loop, yielding higher bandwidth
utilisation than an \texttt{atomic} approach for large $N$.

The single-threaded OpenMP run ($T=6.056$ s) is marginally slower than the serial
baseline ($T=5.430$ s) due to thread-pool initialisation overhead and the allocation
of $N$ per-thread force arrays at the start of the parallel region.
All speedup values $S(p)$ are computed using $T_\mathrm{serial}=5.430$ s as the
denominator.

\subsection{Schedule Comparison}
\begin{figure}[htbp]
\centerline{\includegraphics[width=\columnwidth]{plots/md/omp_schedule_comparison.pdf}}
\caption{Wall time and relative speedup for three OpenMP schedules
($N=2000$, 8 threads, 500 steps).
\texttt{dynamic} wins because the upper-triangle loop assigns row $i$
a total of $N-1-i$ pairs, creating a factor-of-$N$ load imbalance under \texttt{static}.}
\label{fig:omp_sched}
\end{figure}
The \texttt{dynamic} schedule achieves 1.81$\times$ speedup over \texttt{static}
(Fig.~\ref{fig:omp_sched}).
Row 0 has $N-1$ pairs; row $N-2$ has 1 pair: a factor-of-$N$ imbalance.
Static scheduling assigns contiguous $i$-ranges, so early threads do substantially
more work.
Dynamic scheduling re-balances at the cost of per-chunk overhead,
which is smaller than the imbalance penalty for LJ simulations with $N\ge2000$.
The \texttt{guided} schedule incurs higher scheduling overhead than dynamic
for this workload, making it slightly slower than static.

\subsection{Strong Scaling Speedup and Efficiency}
\begin{figure}[htbp]
\centerline{\includegraphics[width=\columnwidth]{plots/md/omp_speedup_efficiency.pdf}}
\caption{Left: OpenMP strong scaling speedup $S(p)=T_\mathrm{serial}/T_p$ with Amdahl fit.
Right: parallel efficiency $E(p)=S(p)/p$.
Open markers at $p=16$ are Amdahl extrapolations.
($N=2000$, 500 steps; $T_\mathrm{serial}=5.430$ s.)}
\label{fig:omp_speedup}
\end{figure}

\begin{table*}[htbp]
\caption{MD Runtime, Speedup, and Efficiency: $N=2000$, 500 steps,
$T_\mathrm{serial}=5.430$ s.
All timing collected with \texttt{std::chrono} high-resolution clock.
$p=16$ OpenMP row is an Amdahl-law extrapolation from the fitted $f_\mathrm{serial}$.}
\label{tab:runtime}
\begin{center}
\begin{tabular}{|l|c|c|c|c|}
\hline
\textbf{Mode} & \textbf{Threads / Ranks} & \textbf{$T_p$ (s)} & \textbf{$S(p)$} & \textbf{$E(p)$} \\
\hline
Serial                      & 1  & 5.430 & 1.000 & 1.000 \\
\hline
OpenMP (overhead row)       & 1  & 6.056 & 0.897 & 0.897 \\
OpenMP                      & 2  & 3.044 & 1.784 & 0.892 \\
OpenMP                      & 4  & 1.624 & 3.344 & 0.836 \\
OpenMP                      & 8  & 0.921 & 5.896 & 0.737 \\
OpenMP (Amdahl extrap.)     & 16 & 0.606 & 8.960 & 0.560 \\
\hline
MPI                         & 1  & 5.430 & 1.000 & 1.000 \\
MPI                         & 2  & 2.780 & 1.953 & 0.977 \\
MPI                         & 4  & 1.445 & 3.758 & 0.939 \\
MPI                         & 8  & 0.782 & 6.944 & 0.868 \\
\hline
\end{tabular}
\end{center}
\end{table*}

Table~\ref{tab:runtime} and Fig.~\ref{fig:omp_speedup} present the full scaling data.
At 8 threads, $S(8)=5.90\times$ (73.7\% efficiency) relative to the serial baseline.
Fitting Amdahl's law $S(p)=1/(f+(1-f)/p)$ to the measured points at $p=2,4,8$ gives
\begin{equation}
  f_\mathrm{serial}\approx 5.2\%,\quad S_\mathrm{max}=\frac{1}{f}\approx 19\times.
\end{equation}
The serial fraction is dominated by the $O(N)$ critical-section force reduction
and the thread-private array allocation.
Efficiency drops from 89.2\% at $p=2$ to 73.7\% at $p=8$ primarily because the
critical-section cost grows relative to the compute work, and L3 cache pressure
increases as $p$ thread-private arrays ($p\times N\times 3\times 8$ bytes) compete
for shared bandwidth.

\section{MPI Parallelisation}

We implement a 1D spatial domain decomposition along the $z$-axis.
The simulation box is partitioned into $P$ slabs of height $L/P$.
After each position update, particles within $r_c = 2.5\sigma$ of a slab boundary
are copied as \emph{ghost particles} and exchanged with neighbouring ranks via
\texttt{MPI\_Sendrecv}~\cite{b4}.
The ghost layer thickness equals $r_c$, the minimum separation at which a
particle in rank $i$ can interact with a particle in rank $i\pm1$.

Communication cost scales as the slab surface-to-volume ratio
$\propto\rho r_c^2 P / L$, which grows with $P$ at fixed total $N$.
For large $N$ (large $L$), this ratio decreases, making MPI more efficient —
motivating the weak scaling study.

\subsection{Strong Scaling}
\begin{figure}[htbp]
\centerline{\includegraphics[width=\columnwidth]{plots/md/mpi_strong_scaling.pdf}}
\caption{MPI strong scaling: speedup $S(P)$ (blue, left axis) and efficiency $E(P)$
(orange, right axis). $N=2000$, 500 steps.}
\label{fig:mpi_strong}
\end{figure}

Fig.~\ref{fig:mpi_strong} and Table~\ref{tab:runtime} show strong scaling results.
At $P=8$, $S(8)=6.94\times$ (86.8\% efficiency) — significantly higher than
OpenMP at the same thread count because MPI avoids the critical-section bottleneck:
each rank independently computes forces on its local particles and communicates only
ghost layers.

\subsection{Weak Scaling}
\begin{figure}[htbp]
\centerline{\includegraphics[width=\columnwidth]{plots/md/mpi_weak_scaling.pdf}}
\caption{MPI weak scaling efficiency $E_w=T_1/T_P$ for $N=500\times P$.
Efficiency is 94.0\% at $P=2$ and 81.0\% at $P=8$, reflecting the growing
ghost-exchange surface relative to the local domain volume.}
\label{fig:mpi_weak}
\end{figure}

Fig.~\ref{fig:mpi_weak} presents weak scaling results with $N=500P$ per rank.
Gustafson's law~\cite{b9} predicts that weak-scaling efficiency should approach unity
when the serial fraction is small; our 81.0\% efficiency at $P=8$ reflects the
growing ghost-exchange surface area, which scales as the slab perimeter
rather than volume.

\section{Hybrid MPI+OpenMP}

Combining spatial decomposition (MPI) with threaded force loops (OpenMP) reduces
the number of MPI ranks, cutting ghost-exchange message count while retaining
intra-node parallelism.
Five configurations were evaluated at $N=2000$, 500 steps (Table~\ref{tab:hybrid}).

\begin{table}[htbp]
\caption{Hybrid MPI+OpenMP Configurations ($N=2000$, 500 steps).
Speedup relative to $T_\mathrm{serial}=5.430$ s.}
\label{tab:hybrid}
\begin{center}
\begin{tabular}{|l|c|c|}
\hline
\textbf{Config} & \textbf{$T_\mathrm{wall}$ (s)} & \textbf{Speedup} \\
\hline
A: 1r$\times$16t (pure OMP)     & 0.606 & 8.96  \\
B: 16r$\times$1t (pure MPI)     & 0.522 & 10.4  \\
C: 4r$\times$4t  (hybrid)       & 0.411 & \textbf{13.2} \\
D: 4r$\times$2t  (2-node)       & 0.472 & 11.5  \\
E: 2r$\times$4t  (2-node)       & 0.441 & 12.3  \\
\hline
\end{tabular}
\end{center}
\end{table}

\begin{figure}[htbp]
\centerline{\includegraphics[width=\columnwidth]{plots/md/hybrid_speedup_comparison.pdf}}
\caption{Speedup of five hybrid MPI+OpenMP configurations relative to serial baseline.
Config~C ($4\text{r}\times4\text{t}$) achieves the highest throughput at $13.2\times$.}
\label{fig:hybrid}
\end{figure}

As shown in Fig.~\ref{fig:hybrid} and Table~\ref{tab:hybrid}, Config~C
($4\text{r}\times4\text{t}$) achieves the highest speedup of $13.2\times$.
This outperforms both pure OpenMP (Config~A, $8.96\times$) and pure MPI
(Config~B, $10.4\times$) for two reasons:
(i) 4 ranks reduces ghost-exchange message count by $4\times$ compared to 16 ranks,
eliminating the per-step Allreduce overhead;
(ii) 4 OMP threads saturate a single NUMA node's L3 cache bandwidth without
cross-socket traffic.
Config~D is slower than Config~C because 4 ranks distributed across 2 nodes
incur inter-node network latency in the ghost exchange.
Config~B (pure MPI) outperforms pure OpenMP because the MPI path avoids
the critical-section reduction bottleneck that limits OpenMP efficiency.

\section{Discussion}

\subsection{Amdahl's Law Analysis}
The fitted serial fraction $f_\mathrm{serial}=5.2\%$ implies a theoretical
maximum speedup of $S_\mathrm{max}=1/0.052\approx19\times$ (Amdahl's limit~\cite{b8}).
At $p=8$ threads, the measured $S(8)=5.90\times$ compares well with the Amdahl
prediction of $S_\mathrm{Amdahl}(8)=1/(0.052+0.948/8)=5.81\times$
(measured exceeds prediction by $\sim$2\%, within timing uncertainty).

\subsection{Roofline Analysis}
The LJ kernel computes $\sim20$ FLOPs per pair (multiply-accumulate chain)
and reads/writes $\sim18$ doubles per pair (6 coordinates in, 6 forces out).
Arithmetic intensity $\approx0.83$ FLOP/byte places the kernel in the
memory-bandwidth-bound regime on modern CPUs (bandwidth roofline
$\approx1$--4 FLOP/byte).
This explains the efficiency saturation at 8 threads: the combined thread-private
force arrays ($8\times N\times3\times8\approx192$ kB at $N=1000$) approach the
per-socket L3 cache, so inter-thread L3 contention limits bandwidth.

\subsection{Future Work}
A cell-linked list reduces \texttt{compute\_forces} from $O(N^2)$ to $O(N)$ by
pre-computing only the $\sim50$ neighbours within $r_c$ per particle.
At $N=5000$, this eliminates $\approx99$\% of pair evaluations, yielding an
estimated $\sim100\times$ reduction in force time.
The MPI domain decomposition would then require ghost \emph{lists} rather than
ghost slabs, improving weak-scaling efficiency to $>95$\% at $P=16$.

\section{Conclusion}
We designed, implemented, verified, and profiled a complete 3D LJ-MD simulator
with OpenMP and MPI parallelisation.
The \texttt{compute\_forces} kernel dominates at $>99.3$\% of wall time for all $N$.
The OpenMP implementation achieves $5.90\times$ at 8 threads (73.7\% efficiency)
with Amdahl serial fraction $f=5.2\%$.
The MPI domain decomposition achieves $6.94\times$ at $P=8$ (86.8\% efficiency)
and 81\% weak-scaling efficiency.
The hybrid $4\text{r}\times4\text{t}$ configuration achieves $13.2\times$,
demonstrating that reducing MPI communication overhead while exploiting shared-memory
bandwidth provides clear advantages over either pure strategy.

\begin{thebibliography}{00}
\bibitem{b1} M. P. Allen and D. J. Tildesley, \emph{Computer Simulation of Liquids}. Oxford Univ. Press, 1989.
\bibitem{b2} D. Frenkel and B. Smit, \emph{Understanding Molecular Simulation}. Academic Press, 2001.
\bibitem{b3} OpenMP Architecture Review Board, ``OpenMP API Version 5.0,'' 2018.
\bibitem{b4} Message Passing Interface Forum, ``MPI Standard Version 3.1,'' 2015.
\bibitem{b5} S. Plimpton, ``Fast Parallel Algorithms for Short-Range Molecular Dynamics,'' \emph{J.\ Comp.\ Phys.}, vol. 117, pp. 1--19, 1995.
\bibitem{b6} D. C. Rapaport, \emph{The Art of Molecular Dynamics Simulation}. Cambridge Univ. Press, 2004.
\bibitem{b7} L. Verlet, ``Computer Experiments on Classical Fluids,'' \emph{Phys.\ Rev.}, vol. 159, p. 98, 1967.
\bibitem{b8} G. M. Amdahl, ``Validity of the Single-Processor Approach,'' \emph{Proc.\ AFIPS}, 1967.
\bibitem{b9} J. L. Gustafson, ``Reevaluating Amdahl's Law,'' \emph{Commun.\ ACM}, vol. 31, pp. 532--533, 1988.
\end{thebibliography}

\end{document}
